\section*{附录：各问核心程序}
\addcontentsline{toc}{section}{附录：各问核心程序}

本次训练赛仅提交 PDF，因此附录直接给出四问正式求解程序的核心代码节选。为控制篇幅，数据读取、绘图、文件导出和重复审计代码不再展开；下列代码分别对应正文中的 MCP 基础调度、SFETA、BESS-ED 与 Benders--列生成联合求解。

\begingroup
\lstset{language=Python,basicstyle=\ttfamily\scriptsize,columns=fullflexible,breaklines=true,showstringspaces=false}

\subsection*{A.1 问题一：MCP 概率预测与本地优先基础调度}
\begin{lstlisting}[caption={问题一核心代码节选（q1\_solver.py）}]
def compound(rate, vals, n=5000):
    ns = rng.poisson(rate, size=n)
    ans = np.zeros(n, dtype=float)
    for j, m in enumerate(ns):
        if m:
            ans[j] = rng.choice(vals, size=m, replace=True).sum()
    return ans

refit = work.loc[work.ArrivalHour <= 2375]
lam2 = len(refit) / 2376.0
pi2 = refit.groupby(["SourceRegion", "TaskType"]).size() / len(refit)
em2 = marks(refit)
for r in regions:
    for k in types:
        rate = lam2 * float(pi2.get((r, k), 0.0))
        sim = compound(rate, em2[k][:, 0])
        vals2["PI05"][(r, k)] = float(np.quantile(sim, 0.05))
        vals2["PI95"][(r, k)] = float(np.quantile(sim, 0.95))

# 对测试窗实际任务执行“本地完整时间域优先，再考虑远端”
chosen = find_candidate([src])
if chosen is None:
    other = sorted(
        [r for r in regions if r != src and lmap[(src, r)] <= x.MaxLatency_ms],
        key=lambda r: (lmap[(src, r)], r),
    )
    chosen = find_candidate(other)
if chosen is None:
    bad.append(int(x.TaskID))
else:
    r, st = chosen
    update_occupancy(r, st, x)
\end{lstlisting}

\subsection*{A.2 问题二：SFETA 完整合法域动态边际任务分配}
\begin{lstlisting}[caption={问题二核心代码节选（q2\_iterative\_solver.py）}]
def build_priority(tasks, latency, region_index):
    work = tasks.copy()
    work["RegionCount"] = work.apply(
        lambda row: int(np.sum(
            latency[region_index[row.SourceRegion], :] <= row.MaxLatency_ms + 1e-9)),
        axis=1,
    )
    work["Slack_h"] = (
        work.LatestFinishHour - work.ArrivalHour
        - work.EstimatedDuration_min / 60.0
    )
    work["GPUHour"] = work.GPU_Demand * work.EstimatedDuration_min / 60.0
    return work.sort_values(
        ["RegionCount", "Slack_h", "GPUHour", "TaskID"],
        ascending=[True, True, False, True],
    )

for task in work.itertuples(index=False):
    core.add_task(gpu, ai, task, old_region, old_start,
                  power, profile_cache, -1.0)
    for region in allowed:
        starts = np.arange(lo, hi + 1, dtype=int)
        valid, dcurt, dcost, dcarbon = core.evaluate_region_candidates(
            int(region), starts, profile, float(task.GPU_Demand),
            float(task.GPU_Demand) * float(power[task.TaskType]),
            gpu, ai, base_facility, nonai, capacity_gpu, max_it,
            pue, max_facility, c0, g0, grid0, used0,
            price, carbon, max_grid,
        )
        collect_feasible_candidates(region, starts, valid,
                                    dcurt, dcost, dcarbon)
    choice = core.pick_candidate(all_candidates, mode)
    core.add_task(gpu, ai, task, choice.region, choice.start,
                  power, profile_cache, 1.0)
\end{lstlisting}

\subsection*{A.3 问题三：BESS-ED 两层字典序线性规划}
\begin{lstlisting}[caption={问题三核心代码节选（q3\_solver.py）}]
# 第一层：净运行成本最小化
c_cost = np.zeros(n)
c_cost[off["gL"]:off["gL"] + T] = buy
c_cost[off["qG"]:off["qG"] + T] = buy
c_cost[off["s"]:off["s"] + T] = -sell

r1 = linprog(
    c_cost, A_ub=Aub_base, b_ub=bub_base,
    A_eq=Aeq, b_eq=beq, bounds=bnds, method="highs"
)
if not r1.success:
    raise RuntimeError(r1.message)
cost_star = float(r1.fun)

# 第二层：在成本最优面内最小化储能吞吐量
c_throughput = np.zeros(n)
for name in ("qR", "qG", "d"):
    c_throughput[off[name]:off[name] + T] = 1.0
cost_row = coo_matrix(c_cost.reshape(1, -1))
A2 = vstack([Aub_base, cost_row], format="csr")
b2 = np.r_[bub_base, cost_star + TOL_COST]
r2 = linprog(
    c_throughput, A_ub=A2, b_ub=b2,
    A_eq=Aeq, b_eq=beq, bounds=bnds, method="highs"
)
if not r2.success:
    raise RuntimeError(r2.message)
\end{lstlisting}

\subsection*{A.4 问题四：区域 Benders 与完整域列生成}
\begin{lstlisting}[caption={问题四核心代码节选（q4\_full\_solver.py）}]
def price_all(pool, data, gpu_row, it_row,
              rdual, cdual, cuts, eqdual, tol):
    """所有合法 task-region-start 均参与精确定价。"""
    additions = []
    for i in range(len(data.task)):
        start0 = int(data.arrival[i])
        end = (start0 if data.task.iloc[i].TaskType == "RealTimeInference"
               else legal_last_start(data, i))
        best = None
        for r in np.flatnonzero(data.legal[i]):
            scores = reduced_cost_profile(i, r, start0, end,
                                          rdual, cdual, cuts)
            s = first_inactive_argmin(pool, i, r, start0, scores)
            if s is not None:
                value = float(scores[s - start0] * data.gpu[i])
                if best is None or value < best[0]:
                    best = (value, int(r), int(s))
        if best is not None:
            reduced = best[0] - eqdual[i]
            if reduced < -tol:
                additions.append((i, best[1], best[2], reduced))
    return sorted(additions, key=lambda z: z[3])

for iteration in range(1, max_iter + 1):
    rmp, _, _, _ = solve_rmp(pool, data, gpu_row, it_row, rhs, cuts)
    load = facility_load(pool, data, rmp["x"])
    energy = energy_lp(load, data)
    violation = float(energy["cost"] - rmp["theta"])
    if violation > benders_tol:
        cuts.append(make_region_cut(energy, load, data))
        continue
    additions = price_all(pool, data, gpu_row, it_row,
                          rdual, cdual, cuts, rmp["eq_dual"], price_tol)
    if additions:
        for i, r, s, _ in additions:
            pool.add(i, r, s)
        continue
    root_closed = True
    break
\end{lstlisting}

\endgroup

问题四在 Cost 层闭合后继续对 Wait、Latency 两层执行相同的完整域定价与真实能源复核，并在最终扩展列池上进行两轮再认证。新能源下降 20\% 场景只保留严格 Cost 上下界及经硬约束审计的可行候选，不将其候选服务质量指标表述为全局最优。
